\section{Motivation und Problemstellung}
Obwohl die knotenbasierten Detektionsmechanismen der betrachteten Ladeinfrastruktur bereits eine beachtliche Erkennungsleistung aufweisen (siehe Abbildung~\ref{fig:auszug_threat_report}), besteht das Ziel, die Spezifität (True-Negative-Rate) sowie die Robustheit des Gesamtsystems weiter zu steigern. Lokale Anomalieerkennungsverfahren stoßen systembedingt an Grenzen, wenn netzwerkweite Veränderungen oder subtile, koordinierte Manipulationsversuche vorliegen \cite{lo2022egraphsage}, die auf Ebene einer einzelnen Wallbox nicht als bösartig identifiziert werden können.

Die Motivation dieser Arbeit liegt daher in der Entwicklung einer Methodik, die es einzelnen Ladepunkten ermöglicht, Kontextinformationen benachbarter Knoten in den Entscheidungsprozess einzubeziehen. Es wird vermutet, dass durch diesen kollaborativen Ansatz räumliche und zeitliche Korrelationen innerhalb des Netzwerks besser erfasst werden können, die über die isolierte Betrachtung eines einzelnen Knotens hinausgehen. Durch die Integration dieser Kontextinformationen wird erhofft, eine präzisere Differenzierung zwischen globalen Schwankungen und lokalen Sicherheitsvorfällen zu ermöglichen. Es wird erwartet, dass dieser erweiterte Kontextbezug zu einer Reduktion von Fehlalarmen (False Positives) sowie einer verbesserten Erkennungsrate tatsächlicher Bedrohungen (True Positives) beiträgt. Inwiefern diese Annahmen empirisch belastbar sind, bleibt jedoch offen und soll in einer weiterführenden Arbeit durch die Evaluation anhand eines unabhängigen Datensatzes systematisch untersucht werden.